% ---------------------------------------------------------------------------
% 2 DESENVOLVIMENTO
% ---------------------------------------------------------------------------
\chapter{Desenvolvimento}\label{cap:desenvolvimento}

% ---------------------------------------------------------------------------
% 2.1 Objetivos
% ---------------------------------------------------------------------------
\section{Objetivos}\label{sec:objetivos}

O objetivo geral define o que se pretende atingir com o projeto.

Os objetivos específicos definem as etapas do trabalho a serem realizadas
para que se alcance o objetivo geral. Os objetivos podem ser: exploratórios,
descritivos e explicativos.

Utilize verbos nos infinitivos para os objetivos:

\begin{itemize}
    \item Exploratórios (conhecer, identificar, levantar, descobrir);
    \item Descritivos (caracterizar, descrever, traçar, determinar);
    \item Explicativos (analisar, avaliar, verificar, explicar).
\end{itemize}

% ---------------------------------------------------------------------------
% 2.2 Justificativa e delimitação do problema
% ---------------------------------------------------------------------------
\section{Justificativa e delimitação do problema}\label{sec:justificativa}

Para a formulação do problema, o grupo deve elaborar uma pergunta que
norteará o desenvolvimento da pesquisa e para a qual será gerada a solução.

Neste item, espera-se que o grupo traga as razões ou práticas que
justifiquem a proposta inicial. Exemplos de justificativa:

\begin{itemize}
    \item Relevância social, cultural e acadêmica;
    \item Contribuições da pesquisa para o local onde o projeto será
          desenvolvido.
\end{itemize}

% ---------------------------------------------------------------------------
% 2.3 Fundamentação teórica
% ---------------------------------------------------------------------------
\section{Fundamentação teórica}\label{sec:fundamentacao}

Pesquisar em fontes confiáveis como monografias, trabalhos de conclusão de
cursos, artigos científicos, revistas especializadas, dissertações e teses,
entre outras fontes, como instituições públicas ligadas às normatizações.

A fundamentação deve ser condizente com o problema em estudo.

Busque e cite fundamentos relevantes e atuais sobre o assunto a ser estudado
e demonstre o entendimento da literatura existente sobre o tema.

As citações e paráfrases devem ser feitas de acordo com as regras da ABNT
6023, de 2002.

Para citações indiretas: (AUTOR, ano) ou Autor (ano).

Para citações diretas:

\begin{itemize}
    \item Menos de três linhas: entre aspas, acompanhadas de (AUTOR, ano,
          p.~xx).
    \item Mais de três linhas: sem aspas, fonte tamanho 10, e recuo de
          parágrafo de 4~cm e espaçamento simples. Exemplo:
\end{itemize}

\begin{citacao}
    Faz necessária a busca por alternativas para dinamizar o processo de
    ensino-aprendizagem em que o professor e os alunos sejam sujeitos e
    caminhem juntos na aventura de aprender e descobrir o novo e vejam
    sentido nos seus fazeres e não simplesmente no cumprimento de mais uma
    tarefa. A matemática, portanto, faz parte da vida e pode ser aprendida
    de uma maneira dinâmica, desafiante e divertida. (PILETTI, 1998, p.~102).
\end{citacao}

% ---------------------------------------------------------------------------
% 2.4 Metodologia
% ---------------------------------------------------------------------------
\section{Metodologia}\label{sec:metodologia}

Metodologia refere-se aos métodos e instrumentos adotados para a execução do
projeto. Nesta seção, espera-se que o grupo descreva os passos e as
estratégias adotadas para o desenvolvimento do Projeto Integrador.

Assim, indique as estratégias adotadas em cada etapa do projeto:

\noindent -- \textbf{Ouvir e interpretar o contexto:}

\begin{itemize}
    \item Descrição do contexto em que o projeto foi realizado;
    \item Perfil dos sujeitos participantes, se for o caso;
    \item Como as informações iniciais foram coletadas: observação,
          entrevista, formulário, questionário etc.
\end{itemize}

\noindent -- \textbf{Criar / Prototipar:}

\begin{itemize}
    \item Análise dos dados, por exemplo, estratégias referentes à pesquisa
          qualitativa ou quantitativa;
    \item Descrição das soluções encontradas ou desenvolvidas para o problema
          investigado.
\end{itemize}

\noindent -- \textbf{Implementar / Testar:}

\begin{itemize}
    \item Como a solução foi testada? Que devolutivas sobre a solução o
          grupo conseguiu coletar?
    \item Que melhorias foram indicadas para as soluções
          propostas/desenvolvidas?
\end{itemize}

Finalmente, este é o espaço para que o leitor do seu projeto entenda, em
detalhes, quais foram as estratégias usadas para que os resultados fossem
obtidos.

% ---------------------------------------------------------------------------
% 2.5 Resultados preliminares: solução inicial
% ---------------------------------------------------------------------------
\section{Resultados preliminares: solução inicial}\label{sec:resultados}

O grupo deve demonstrar a criação de soluções com base na metodologia
indicada pela UNIVESP, respeitando os passos \textbf{ouvir, criar} e
\textbf{implementar}. Portanto, deve identificar quais foram os resultados
obtidos em cada um dos passos para a construção da solução.

É importante que o grupo inclua imagens, \textit{storyboards} ou
ilustrações que demonstrem visualmente a solução adotada, junto aos passos
desenvolvidos. Dessa forma, sugere-se que, neste capítulo, seja apresentada
uma descrição detalhada de como se deu o processo de construção da primeira
solução desenvolvida pelo grupo.

\textbf{Importante}: quando se tratar de projetos desenvolvidos com a
participação de crianças e adolescentes, não é permitida a inclusão de fotos
deles sem a autorização de seus pais ou responsáveis.
